\documentclass{SHVpaper}
\begin{document}
% Так позначається початок рядка із коментарем. Він не компілюється. Далі по тексту описуються особливості використання стилю тез
% Просимо уникати використання додаткових стильових пакетів
% Просимо НЕ використовувати власні макроси чи команди

%

\section{WebRTC конференції без медіасервера: підхід на основі обрання вузлів ретрансляторів}

\authors{Романків Святослав Степанович}

%Перед відправкою тез, будь ласка, переконайтесь у відсутності орфографічних та змістовних помилок, та скомпілюйте документ в pdf.
%На електронну пошту надішліть, будь ласка, tex-файл та pdf-файл.
%Далі наводимо короткий огляд основних результатів, які планується представити в доповіді
%Обсяг скомпільованих тез НЕ повинен перевищувати одну сторінку

Web Real-Time Communication (WebRTC) --- відкритий протокол, основною метою якого є забезпечення обміну аудіо, відео та іншими довільними даними між учасниками через найбільш прямий шлях з’єднання \cite{rfc8825}. Підтримка наявності більше двох учасників у конференції традиційно вимагає серверної медіаінфраструктури. У SFU архітектурі медіасервер пересилає потоки між учасниками без модифікацій, у MCU --- змішує їх на стороні сервера. Обидва підходи потребують ресурсів, пропорційних кількості учасників у сесії. Розглянутою альтернативою є обрання самих учасників у якості ретрансляторів медіапотоків.

Кожен учасник вимірює декілька основних метрик: пропускну здатність, тип NAT, кругову затримку, стабільність з’єднання та можливості пристрою. Метрики транслюються через канал сигналізації.

Зважена функція оцінювання для кожного учасника має вигляд:

\begin{equation}
\label{score_formula}
s = w_b \cdot S_b + w_n \cdot S_n + w_l \cdot S_l + w_s \cdot S_s + w_d \cdot S_d
\end{equation}

Компоненти у (\ref{score_formula}) нормовані до $[0,1]$. Лідер конференції (учасник із найменшим ідентифікатором) формує топологію конференції та транслює її через сервер сигналізації іншим активним користувачам. Після отримання учасником повідомлення про його роль ретранслятора, він активує механізм пересилання та налаштовує таблиці маршрутизації медіа.

Переоцінка відбувається періодично та у разі суттєвих змін метрик. Ретранслятор, оцінка якого опускається нижче найнижчої оцінки серед звичайних учасників, понижується в ролі та замінюється.

Запропонований підхід дозволяє організувати багатосторонні WebRTC конференції без виділених медіасерверів, що дозволяє знизити операційні витрати та усунути централізовані точки відмови. Функцію ретрансляції виконують самі учасники, обрані на основі їхніх мережевих характеристик.

\begin{thebibliography}{5}
%в дужках вказуэться ім'я, яке використовується для генерації посилань в тексті
\bibitem{rfc8825} Alvestrand H. Overview: Real-Time Protocols for Browser-Based Applications. -- IETF RFC 8825, https://datatracker.ietf.org/doc/html/rfc8825
\end{thebibliography}

\subsection{Автори}
%для кожного автора потрібно вказати детальну інформацію: ПІБ, наукове звання та місце навчання, електронну адресу
\author{Романків Святослав Степанович}{студент 2-го курсу магістратури, факультет комп'ютерних наук та кібернетики, Київський національний університет імені Тараса Шевченка, Київ, Україна}{romankivs@knu.com}

\end{document}
